% concatenated from main.tex
\documentclass[preprint]{imsart}

\RequirePackage{fix-cm}






\usepackage{latexsym}
\usepackage{array}
\usepackage{color}
\usepackage{rotating}
\usepackage{epsfig, wrapfig}
\usepackage{graphicx, lscape, graphics, color} % removed subfig 
\usepackage{verbatim}
\usepackage{latexsym}
\usepackage{fancyvrb}
\usepackage{pdfpages}
\usepackage{fmtcount}
\usepackage{multirow}
\usepackage{url}
\usepackage[english]{babel}
\usepackage[utf8]{inputenc}
\usepackage{tikz} 
\usetikzlibrary{arrows, decorations.pathmorphing, backgrounds, fit, positioning, shapes.symbols, chains}
\usetikzlibrary{decorations.markings}
\usepackage{lscape}
\usepackage{epstopdf}
\usepackage{algpseudocode}
\usepackage{algorithm}
\usepackage{kotex}
\usepackage{soul}
\usepackage{enumerate}
\usepackage{colortbl}
\usepackage{booktabs}
\usepackage{appendix}
\usepackage{mathtools}
\usepackage{tikz} 
\usetikzlibrary{arrows, decorations.pathmorphing, backgrounds, fit, positioning, shapes.symbols, chains}
\usetikzlibrary{decorations.markings}
\usetikzlibrary{arrows.meta}


\newcommand\norm[1]{\left\lVert#1\right\rVert}
\newcommand{\eighth}{{\textstyle \frac{1}{8}}}
%%%%%% Operator
\DeclareMathOperator*{\argmin}{arg\,min}
\DeclareMathOperator*{\argmax}{arg\,max}
\DeclarePairedDelimiter\ceil{\lceil}{\rceil}
\DeclarePairedDelimiter\floor{\lfloor}{\rfloor}
\newcommand{\Tau}{\mathcal{T}}
\newcommand{\cA}{\mathcal{A}}
\newcommand{\cB}{\mathcal{B}}
\newcommand{\cC}{\mathcal{C}}
\newcommand{\cD}{\mathcal{D}}
\newcommand{\cE}{\mathcal{E}}
\newcommand{\cF}{\mathcal{F}}
\newcommand{\cG}{\mathcal{G}}
\newcommand{\cH}{\mathcal{H}}
\newcommand{\cI}{\mathcal{I}}
\newcommand{\cJ}{\mathcal{J}}
\newcommand{\cK}{\mathcal{K}}
\newcommand{\cL}{\mathcal{L}}
\newcommand{\cM}{\mathcal{M}}
\newcommand{\cN}{\mathcal{N}}
\newcommand{\cO}{\mathcal{O}}
\newcommand{\cP}{\mathcal{P}}
\newcommand{\cQ}{\mathcal{Q}}
\newcommand{\cR}{\mathcal{R}}
\newcommand{\cS}{\mathcal{S}}
\newcommand{\cT}{\mathcal{T}}
\newcommand{\cU}{\mathcal{U}}
\newcommand{\cV}{\mathcal{V}}
\newcommand{\cW}{\mathcal{W}}
\newcommand{\cX}{\mathcal{X}}
\newcommand{\cY}{\mathcal{Y}}
\newcommand{\cZ}{\mathcal{Z}}
\newcommand{\bbA}{\mathbb{A}}
\newcommand{\bbB}{\mathbb{B}}
\newcommand{\bbC}{\mathbb{C}}
\newcommand{\bbD}{\mathbb{D}}
\newcommand{\bbE}{\mathbb{E}}
\newcommand{\bbF}{\mathbb{F}}
\newcommand{\bbG}{\mathbb{G}}
\newcommand{\bbH}{\mathbb{H}}
\newcommand{\bbI}{\mathbb{I}}
\newcommand{\bbJ}{\mathbb{J}}
\newcommand{\bbK}{\mathbb{K}}
\newcommand{\bbL}{\mathbb{L}}
\newcommand{\bbM}{\mathbb{M}}
\newcommand{\bbN}{\mathbb{N}}
\newcommand{\bbO}{\mathbb{O}}
\newcommand{\bbP}{\mathbb{P}}
\newcommand{\bbQ}{\mathbb{Q}}
\newcommand{\bbR}{\mathbb{R}}
\newcommand{\bbS}{\mathbb{S}}
\newcommand{\bbT}{\mathbb{T}}
\newcommand{\bbU}{\mathbb{U}}
\newcommand{\bbV}{\mathbb{V}}
\newcommand{\bbW}{\mathbb{W}}
\newcommand{\bbX}{\mathbb{X}}
\newcommand{\bbY}{\mathbb{Y}}
\newcommand{\bbZ}{\mathbb{Z}}
\newcommand{\bA}{\mathbf{A}}
\newcommand{\bB}{\mathbf{B}}
\newcommand{\bC}{\mathbf{C}}
\newcommand{\bD}{\mathbf{D}}
\newcommand{\bE}{\mathbf{E}}
\newcommand{\bF}{\mathbf{F}}
\newcommand{\bG}{\mathbf{G}}
\newcommand{\bH}{\mathbf{H}}
\newcommand{\bI}{\mathbf{I}}
\newcommand{\bJ}{\mathbf{J}}
\newcommand{\bK}{\mathbf{K}}
\newcommand{\bL}{\mathbf{L}}
\newcommand{\bM}{\mathbf{M}}
\newcommand{\bN}{\mathbf{N}}
\newcommand{\bO}{\mathbf{O}}
\newcommand{\bP}{\mathbf{P}}
\newcommand{\bQ}{\mathbf{Q}}
\newcommand{\bR}{\mathbf{R}}
\newcommand{\bS}{\mathbf{S}}
\newcommand{\bT}{\mathbf{T}}
\newcommand{\bU}{\mathbf{U}}
\newcommand{\bV}{\mathbf{V}}
\newcommand{\bW}{\mathbf{W}}
\newcommand{\bX}{\mathbf{X}}
\newcommand{\bY}{\mathbf{Y}}
\newcommand{\bZ}{\mathbf{Z}}
\newcommand{\fA}{\mathfrak{A}}
\newcommand{\fB}{\mathfrak{B}}
\newcommand{\fC}{\mathfrak{C}}
\newcommand{\fD}{\mathfrak{D}}
\newcommand{\fE}{\mathfrak{E}}
\newcommand{\fF}{\mathfrak{F}}
\newcommand{\fG}{\mathfrak{G}}
\newcommand{\fH}{\mathfrak{H}}
\newcommand{\fI}{\mathfrak{I}}
\newcommand{\fJ}{\mathfrak{J}}
\newcommand{\fK}{\mathfrak{K}}
\newcommand{\fL}{\mathfrak{L}}
\newcommand{\fM}{\mathfrak{M}}
\newcommand{\fN}{\mathfrak{N}}
\newcommand{\fO}{\mathfrak{O}}
\newcommand{\fP}{\mathfrak{P}}
\def\Frechet{Fr\'{e}chet}

\def\av{\mathbf a}
\def\bv{\mathbf b}
\def\cv{\mathbf c}
\def\dv{\mathbf d}
\def\ev{\mathbf e}
\def\fv{\mathbf f}
\def\gv{\mathbf g}
\def\hv{\mathbf h}
\def\iv{\mathbf i}
\def\jv{\mathbf j}
\def\gv{\mathbf g}
\def\kv{\mathbf k}
\def\lv{\mathbf l}
\def\mv{\mathbf m}
\def\nv{\mathbf n}
\def\ov{\mathbf o}
\def\pv{\mathbf p}
\def\qv{\mathbf q}
\def\rv{\mathbf r}
\def\sv{\mathbf s}
\def\tv{\mathbf t}
\def\uv{\mathbf u}
\def\vv{\mathbf v}
\def\wv{\mathbf w}
\def\xv{\mathbf x}
\def\yv{\mathbf y}
\def\zv{\mathbf z}

\RequirePackage{amsthm, amsmath, amsfonts, amssymb}
%\RequirePackage[sort&compress, numbers]{natbib}
\RequirePackage[authoryear]{natbib}%% uncomment this for author-year citations
%\RequirePackage[colorlinks,citecolor=blue,urlcolor=blue]{hyperref}
\usepackage{hyperref}
%\usepackage{cleveref}


%\RequirePackage{graphicx}
%\startlocaldefs
%%%%%%%%%%%%%%%%%%%%%%%%%%%%%%%%%%%%%%%%%%%%%%
%%                                          %%
%% Uncomment next line to change            %%
%% the type of equation numbering           %%
%%                                          %%
%%%%%%%%%%%%%%%%%%%%%%%%%%%%%%%%%%%%%%%%%%%%%%
%\numberwithin{equation}{section}
%%%%%%%%%%%%%%%%%%%%%%%%%%%%%%%%%%%%%%%%%%%%%%
%%                                          %%
%% For Axiom, Claim, Corollary, Hypothesis, %%
%% Lemma, Theorem, Proposition              %%
%% use \theoremstyle{plain}                 %%
%%                                          %%
%%%%%%%%%%%%%%%%%%%%%%%%%%%%%%%%%%%%%%%%%%%%%%
\theoremstyle{plain}
\newtheorem{axiom}{Axiom}
\newtheorem{claim}[axiom]{Claim}
\newtheorem{theorem}{Theorem}
%\newtheorem{lemma}[subsection]{Lemma}
\newtheorem{lemma}{Lemma}
\newtheorem{definition}{Definition}
\newtheorem{remark}{Remark}
\newtheorem{corollary}{Corollary}
\newtheorem{example}{Example}
\newtheorem{proposition}{Proposition}
%%%%%%%%%%%%%%%%%%%%%%%%%%%%%%%%%%%%%%%%%%%%%%
%%                                          %%
%% For Assumption, Definition, Example,     %%
%% Notation, Property, Remark, Fact         %%
%% use \theoremstyle{remark}                %%
%%                                          %%
%%%%%%%%%%%%%%%%%%%%%%%%%%%%%%%%%%%%%%%%%%%%%%
\theoremstyle{remark}

\newtheorem*{fact}{Fact}
%\endlocaldefs
%%%%%%%%%%%%%%%%%%%%%%%%%%%%%%%%%%%%%%%%%%%%%%
%% Please put your definitions here:        %%
%%%%%%%%%%%%%%%%%%%%%%%%%%%%%%%%%%%%%%%%%%%%%%




% \endlocaldefs
%\singlespacing
\begin{document}


\begin{frontmatter}
\title{On the Kolmogorov-Feller weak law of large numbers for the Fr\'{e}chet mean on non-compact symmetric spaces}
%\title{A sample article title with some additional note\thanksref{t1}}
\runauthor{Lee and Jung}
\runtitle{On the Kolmogorov-Feller weak law of large numbers for the Fr\'{e}chet mean}
%\thankstext{T1}{A sample additional note to the title.}

\begin{aug}
%%%%%%%%%%%%%%%%%%%%%%%%%%%%%%%%%%%%%%%%%%%%%%%
%% Only one address is permitted per author. %%
%% Only division, organization and e-mail is %%
%% included in the address.                  %%
%% Additional information can be included in %%
%% the Acknowledgments section if necessary. %%
%% ORCID can be inserted by command:         %%
%% \orcid{0000-0000-0000-0000}               %%
%%%%%%%%%%%%%%%%%%%%%%%%%%%%%%%%%%%%%%%%%%%%%%%
\author[A]{\fnms{Jongmin}~\snm{Lee}\ead[label=e1]{jongmin.lee@pusan.ac.kr}\orcid{0000-0003-1723-4615}}
\and
\author[B]{\fnms{Sungkyu}~\snm{Jung}\ead[label=e2]{sungkyu@snu.ac.kr}\orcid{0000-0002-6023-8956}}
%%%%%%%%%%%%%%%%%%%%%%%%%%%%%%%%%%%%%%%%%%%%%%
%% Addresses                                %%
%%%%%%%%%%%%%%%%%%%%%%%%%%%%%%%%%%%%%%%%%%%%%%
\address[A]{Department of Statistics,
Pusan National University\printead[presep={,\ }]{e1}}

\address[B]{Department of Statistics and Institute for Data Innovation in Science,
Seoul National University\printead[presep={,\ }]{e2}}
\end{aug}

\begin{abstract}
We prove the Kolmogorov–Feller weak law of large numbers for sample Fr\'{e}chet means on non-compact symmetric spaces. The result covers independent, non-identically distributed data, extending beyond the i.i.d. setting. Examples of symmetric positive-definite matrices and product symmetric spaces are provided.
\end{abstract}

\begin{keyword}[class=MSC]
\kwd[Primary ]{62E20}
\kwd{60F05}
%\kwd[; secondary ]{}
\end{keyword}

\begin{keyword}
\kwd{Fr\'{e}chet means}
\kwd{Law of large numbers}
\kwd{Limit theorems}
\kwd{Statistics on manifolds}
\end{keyword}

\end{frontmatter}
%%%%%%%%%%%%%%%%%%%%%%%%%%%%%%%%%%%%%%%%%%%%%%
%% Please use \tableofcontents for articles %%
%% with 50 pages and more                   %%
%%%%%%%%%%%%%%%%%%%%%%%%%%%%%%%%%%%%%%%%%%%%%%
%\tableofcontents
For a sequence of independent and identically distributed real random variables, $(X_i)_{i=1}^{\infty}$, the celebrated Kolmogorov--Feller weak law of large numbers 
\citep[][pp. 234-236]{feller1971introduction} %(see, e.g., \cite{feller1971introduction})
%\citep[KF LLN for short,][pp. 234-236]{feller1971introduction} %\citep{naderi2019version} 
says
 \begin{equation}\label{KF_wlln}
  n\mathbb{P}(|X_1|>n) \underset{n \to \infty}\to 0 ~ \mbox{if and only if} ~ \frac{\sum_{i=1}^{n} X_i - n \mathbb{E}[X_1 1_{|X_1|\le n}]}{n} \overset{\mathbb{P}}{\to} 0 \quad \mbox{as} ~ n \to \infty.
\end{equation}
%{\color{red}The one-dimensional version is classical; see \cite{feller1971introduction}. For completeness, we include a proof of the sufficient direction of (\ref{KF_wlln}) in Section D of the supplementary material.} 
The result in (\ref{KF_wlln}) can be viewed as an extension of the weak law of large numbers for random variables whose distributions are possibly too heavy-tailed to admit a finite first moment (i.e., $\mathbb{E}|X_1| = \infty$), yet still exhibits moderate tail decay in the sense that $n\mathbb{P}(|X_1|>n) \underset{n\to \infty}{\to} 0$, which is called the Kolmogorov--Feller tail condition. 

%%
The Kolmogorov--Feller weak law of large numbers 
has been extended to various settings. For example, it is well-known that \eqref{KF_wlln} holds for an independent but not necessarily identically distributed sequence $(X_i)_{i=1}^{\infty}$ of real-valued random variables \citep[][Theorem 2.2.11]{durrett2019probability}. \cite{stoica2010kolmogorov} investigate the case where $(X_i)_{i=1}^{\infty}$ is assumed to be exchangeable, and a conditional version of the extended Kolmogorov--Feller weak law of large numbers is established by \cite{yuan2015conditional}. 
 Moreover, under a symmetry condition, (\ref{KF_wlln}) 
has been extended for random variables in
%a specific type of 
Banach spaces; see Corollary 3.3 in \cite{marcus1979stable} and Theorem 9.22 in \cite{ledoux1991probability}. However, analogous extensions to nonlinear spaces appear to be unavailable. 


%%
Beyond vector spaces, suppose a data set $\{x_1, x_2, \ldots x_n\}$ lies in a metric space $(M,d)$. For the central tendency of the data set, it is natural to consider the sample \Frechet\ mean, $$\mu_n:=\mbox{argmin}_{m\in M} \frac{1}{n} \sum_{i=1}^{n}d^2(x_i, m),$$ which generalizes the arithmetic mean to any metric space.
Theoretical properties of the sample \Frechet\ mean, such as the law of large numbers, central limit theorem, and non-asymptotic bounds, have been established in Riemannian manifolds and general metric spaces (see, e.g., \cite{ziezold1977expected, sturm2003probability, bhattacharya2003large, schotz2022strong, evans2024limit, brunel2024concentration, schotz2025variance, brunel2025finite}). However, their validity under heavy-tailed distributions that may not possess a finite first moment has not been examined. Moreover, only a few studies to date \citep[e.g.,][]{brunel2024concentration, brunel2025finite} have addressed the case of independent random variables with potentially distinct distributions.


%%
To address this problem, we prove a %KF LLN 
Kolmogorov--Feller-type weak law of large numbers 
for sequences of random variables lying in a specific class of metric spaces. The obstacles for such derivations are two-fold: (i) the possibility of non-uniqueness of \Frechet\ means (for instance, compact manifolds such as the sphere $S^k$ may exhibit non-unique \Frechet\ means), and (ii) the difficulty in evaluating $\mathbb{E}[X_{1}1_{|X_1|\le n}]$ appearing in the right-hand side of (\ref{KF_wlln}). 
We circumvent the non-uniqueness issue by restricting the domain to be a symmetric space of non-compact type (hereafter referred to as a \textit{non-compact symmetric space}), where the sample \Frechet\ mean is known to be unique (e.g., Proposition 4.3 in \cite{sturm2003probability}), albeit the population \Frechet\ mean may not exist. 
%The class of non-compact symmetric spaces encompasses a broad range of metric spaces commonly encountered in statistical data analysis. 
Representative examples include Euclidean space $\mathbb{R}^k$, hyperbolic space $\mathbb{H}^{k}$, the space of symmetric positive-definite matrices $(\text{Sym}^{+}(k), ~ d_{\text{\tiny AI}})$ equipped with the affine-invariant metric, and their product spaces.

The remedy for (ii) is to impose a symmetry condition on the underlying probability distribution $\mathbb{P}_{X}$, as done for the Banach space extension of  \cite{marcus1979stable}.
The symmetry condition, defined precisely in the next section, ensures that there exists a center of symmetry $\mu \in M$ corresponding to $\mathbb{P}_{X}$, and we take $\mu$ as a reference point. To generalize (\ref{KF_wlln}) to our setting, the strategy is to replace $|X_1|$ with $d(X_1, \mu)$, and the sample mean $n^{-1}\sum_{i=1}^{n}X_i$ with the sample \Frechet\ mean $\mu_n$. 

%In this work, 
We establish a modified Kolmogorov--Feller tail condition, which is shown to be sufficient for the convergence of the sample \Frechet\ mean in probability, under the symmetry assumption.  An extension to independent but not necessarily identically distributed random variables is discussed as well. 
%We establish nearly minimal conditions that guarantee the convergence of $\mu_n$ to $\mu$ in probability as the sample size increases. Our framework allows the sequence of random variables to be independent but not necessarily identically distributed. %for i.i.d. random variables  {\color{blue} We further derive a tail condition that ensures the consistency of the sample \Frechet\ mean even in cases where the first moment is infinite, i.e., $\mathbb{E}[d(X, \mu)] = \infty$.}

% %%
% The class of metric spaces known as symmetric spaces of non-compact type (hereafter referred to as non-compact symmetric spaces)
% %(hereafter, referred to as {\it non-compact} symmetric spaces, including the case $M=\mathbb{R}^k$) 
% encompasses a broad class of metric spaces commonly encountered in statistical data analysis. 
% Representative examples of such spaces include Euclidean space $\mathbb{R}^k$, hyperbolic space $\mathbb{H}^{k}$, the space of symmetric positive-definite matrices $(\text{Sym}^{+}(k), ~ d_{\text{\tiny AI}})$ equipped with the affine-invariant metric, and their product spaces.
  
%Our main results provide a generalization of (\ref{KF_wlln}) to this setting. 
%%



%Despite the growing interest in statistics on manifolds, little is known about their theoretical foundations under contamination (e.g., \cite{lin2021functional}). Our results are entirely new to robust statistics on manifolds, and provide sufficient conditions for the weak law of large numbers to hold even under strongly heavy-tailed distributions.


%Little is known about statistical theory and methodology on manifolds in the presence of contamination (e.g., \cite{lin2021functional}). We emphasize that our main results, which are entirely new in the fields of robust statistics and statistics on manifolds, provide sufficient conditions under which the weak law of large numbers holds, even for strongly heavy-tailed distributions.
%%

%%
% The remainder of this article is organized as follows. The next section presents the necessary preliminaries and the main results, accompanied by illustrative examples. Section 3 provides auxiliary results and proofs, some of which are deferred to the supplementary material for better readability.



\section{Kolmogorov--Feller's weak law of large numbers for Fréchet means}

%%
A simply connected and geodesically complete Riemannian manifold is said to be a {\it Hadamard manifold} if it has non-positive curvatures everywhere. A connected, complete and smooth Riemannian manifold $(M,d)$ is said to be a {\it symmetric space} if for any $\nu \in M$, there exists a unique isometry $s_\nu$ (called geodesic symmetry) that fixes $\nu$ and reverses all geodesics through $\nu$; that is, for each $x \in M$, $s_{\nu}(x)=\mbox{Exp}_{\nu}\{-\mbox{Log}_{\nu}(x)\}$, provided it is well-defined. It is known that any non-compact symmetric space is simply connected and has non-positive curvatures everywhere \citep[][]{helgason1979differential}; thus, it is Hadamard. 
% 
%{\color{blue} 
Moreover, for $M$ that is Hadamard, the Riemannian logarithmic map at $p$ for any $p \in M$ is a diffeomorphism between $M$ and $T_{p}M \cong \mathbb{R}^{k}$ by the Cartan-Hadamard theorem. 
%}
 


 
The notion of geodesic symmetry of a probability distribution $\mathbb{P}_{X}$ is defined as follows.
\begin{definition}
Given a symmetric space $(M,d)$ and a probability distribution $\mathbb{P}_{X}$ on the space, if there exists a $\mu \in M$ such that the distribution of $s_\mu(X)$, $\mathbb{P}_{s_\mu (X)}$, equals $\mathbb{P}_X$, then $\mathbb{P}_{X}$ is said to be geodesically symmetric about $\mu$. %{\color{blue}[$s\mu(X)$ is, I believe, is uniquely given. ($s
%_\mu(x) = \exp_\mu(-\log_\mu (x))?$) but in the paragraph above it doesn't say that $s_\mu(x)$ is unique. Without uniqueness, this definition using $s_\mu$ is ambiguous.]}
\end{definition}
A natural question that arises from the definition is whether a given distribution $\mathbb{P}_{X}$ can have two or more points of geodesic symmetry. The answer is negative under the following non-compactness condition.

%A natural question arising from the definition is that for a given distribution $\mathbb{P}_{X}$, if the distribution is geodesically symmetric about $\mu_1, \mu_2 \in M$, then is it true that $\mu_1=\mu_2$? The answer to the question is `yes' under the following non-compactness condition.

% {\color{blue}[I think the finite variance condition in Proposition~\ref{prop:symmetry} can be removed, and it is important to remove the moment condition. 

% A sketch of proof idea: 

% 1. $T = s_{\mu_1} \circ s_{\mu_2}$ is a translation iff $\mu_1\neq \mu_2$. ("transvection", see p.6 of \texttt{https://www-fourier.ujf-grenoble.fr/sites/ifmaquette.ujf-grenoble.fr/files/Maubon.pdf}). Thus $T$, when repeatedly applied, pushes any point to infinity. I.e., for any bounded sets $K$ and $A$, there exists $k$ such that $T^{k'}(A)\cap K = \phi$ for all $k' \ge k$ (think $K$ large, $A$ small). 

% Let $A$ be any bounded Borel set in $M$. Since $P_X$ is geodesically symmetric about both $\mu_1$ and $\mu_2$, $P_X(A) = P_X(T(A))$

% Fix an $\epsilon > 0$. 
% $P_X$ satisfies that there exists a compact set $K = K(\epsilon)$ satisfying $P_X(K) \ge 1-\epsilon$. Thus, for all large $k$ we have $T^k(A) \cap K = \phi$, so that $P(A) = P(T^k(A)) \le 1 - P(K) \le \epsilon$. 

% Sine $\epsilon$ was arbitrary, we have $P(A) = 0$. This is a contradiction and $\mu_1 = \mu_2$ (Details need to be filled)

% ]}





\begin{proposition}\label{prop:symmetry}
Let $(M,d)$ be a non-compact symmetric space. Suppose that $\mathbb{P}_{X}$ is geodesically symmetric about both $\mu_1, \mu_2 \in M$. Then, $\mu_1=\mu_2$.  
\end{proposition}
% \begin{proof}[Proof of Proposition~\ref{prop:symmetry}]
% Suppose, for contradiction, that $\mu_1 \neq \mu_2$. Let $T = s_{\mu_2} \circ s_{\mu_1}$, and $\gamma$ be the geodesic from $\mu_1$ to $\mu_2$. Denote $\ell = 2d(\mu_1, \mu_2) > 0$. Then, note that $T$ is a translation-type isometry of length $\ell$ along the direction of $\gamma$, more precisely, referred to as a transvection. (For an illustration of the transvection, see Figure~\ref{fig:transvection}.) An application of Corollary 2.5 of \cite{sturm2003probability} leads that, since $M$ is Hadamard, we have for any $x \in M$, $d(x, s_{\mu_2}\circ s_{\mu_1}(x))\ge 2d(\mu_1, \mu_2)=\ell$.
% %
% %It can be shown that for any $x \in M$, $d(x, s_{\mu_2}\circ s_{\mu_1}(x))\ge 2d(\mu_1, \mu_2)=\ell$ by an application of Corollary 2.5 of \cite{sturm2003probability}. 
% %
% It implies that for any $x\in M$ and for any $m\in \mathbb{N}$, $d(x, T^m(x))\ge \ell m$. Thus, the isometry $T$, when repeatedly applied, pushes any point to infinity. 
% %%




% %% 
% For given bounded sets $K$ and $A$, define a constant $C_{K, A}=\sup_{k\in K, a\in A}d(k, a) < \infty$. There exists an integer $m$ such that $\ell m> C_{K,A}$. For any $k \in K$ and for any $a \in A$, by triangle inequality, 
% $$
% d(k, T^{m}(a)) \ge d(a, T^{m}(a)) - d(k,a) \ge \ell m - C_{K,A} >0,
% $$
% which implies that $K \cap T^{m}(A) =\phi$, where $T^{m}(A):=\{T^m(a): a \in A\}$. That is, for any bounded sets $K$ and $A$, there exists an integer $m_0$ such that for any integer $m'\ge m_0$, $K \cap T^{m'}(A) = \phi$, where for any $m\ge 1$, $T^{m}(A) \subset M$.


% %%
% We now set the bounded sets $A$ and $K$ as follows. 
% First, $A$ is chosen to be any bounded Borel set in $M$. Since $\mathbb{P}_X$ is geodesically symmetric about both $\mu_1$ and $\mu_2$, it holds that $\mathbb{P}_X(A) = \mathbb{P}_X(T(A))$.

% Notably, any probability measure on $M$ is tight due to the second countability assumption on the manifold $M$ \citep{lee2013smooth}. In words, fix a small $\epsilon > 0$; then, there exists a compact set $K = K(\epsilon)$ satisfying $\mathbb{P}_X(K) \ge 1-\epsilon$. For all large $m$, we have $T^m(A) \cap K = \phi$, so that $\mathbb{P}(A) = \mathbb{P}(T^m(A)) \le 1 - \mathbb{P}(K) \le \epsilon$. 
% %
% Since $\epsilon$ was arbitrary, we have $\mathbb{P}(A) = 0$ for any bounded Borel set $A \subset M$. This is a contradiction to the tightness of $\mathbb{P}_{X}$, and therefore, $\mu_1 = \mu_2$.
% \end{proof}
 




The proof of Proposition~\ref{prop:symmetry} is given in Appendix A. Proposition~\ref{prop:symmetry} implies that any distribution $\mathbb{P}_{X}$ on $M$ admits at most one center of symmetry $\mu$. Geodesically symmetric distributions on $M$ include isotropic distributions as well as elliptically symmetric distributions, which constitute a broad class of probability distributions.
%
The assumption of geodesic symmetry mirrors the symmetry assumption imposed for similar results for Banach spaces \citep{marcus1979stable}. In particular, when $M$ is a vector space, $X$ is geodesically symmetric about $\mu$ if and only if $X -\mu \stackrel{d}{=} -(X - \mu)$.



%%
Unless mentioned otherwise, hereafter we assume that $(M, d)$ is a $k$-dimensional non-compact symmetric space. 
To describe the main results, we assume the following condition. %either of the following conditions. 
%For a fixed $\mu \in M$,
\begin{enumerate}
 \item[(C)] For a fixed $\mu \in M$, $(X_{i})_{i=1}^{\infty}$ is a sequence of independent $M$-valued random variables and for each $i\ge 1$, $\mathbb{P}_{X_{i}}$ is geodesically symmetric about $\mu$.
 %\item[(C2)] $(X_i)_{i=1}^{\infty}$ is a sequence of independent $M$-valued random variables with common distribution $\mathbb{P}_{X}$ that is geodesically symmetric about $\mu$.
\end{enumerate}
Condition (C) imposes only independence on the sequence $(X_i)_{i=1}^{\infty}$, allowing the random variables to have different distributions that are geodesically symmetric about $\mu$.
%Condition (C2), in addition, requires that the random variables share an identical distribution, and thus is a special case of (C1). 



%% 
We denote the sample \Frechet\ mean of $X_1, X_2, \ldots, X_n$ by 
$\mu_n = \mu_{n}(X_1, X_2, \ldots, X_{n})=\mbox{argmin}_{m\in M}n^{-1}\sum_{i=1}^{n} d^2(X_i,m)$.
Since $M$ is a Hadamard manifold, $\mu_n$ is guaranteed to exist and is unique \citep[][Proposition 4.3]{afsari2011riemannian}. On the other hand, under Condition (C), the population \Frechet\ mean, which is the minimizer $p \in M$ of the \Frechet\ function $F(p):= \mathbb{E}[d^2(X_1,p)]$, may not exist, since the second moment $F(p)$ itself may be infinite everywhere. %may not even exist. 
Nevertheless, due to the symmetry about $\mu$ imposed on $(X_i)_{i=1}^{\infty}$, one may intuitively expect that the distribution of $\mu_{n}$ concentrates around the center $\mu$ as $n$ tends to infinity, under suitable tail conditions.  
%
% To state the conditions guaranteeing convergence of $\mu_n$ to $\mu$, we define a triangular array of truncated random variables $(\overline{X_{n, i}})_{n\ge i\ge 1}$ as 
% $$\overline{X_{n, i}} := \begin{cases} X_i \quad \mbox{if} ~ d(X_{i}, \mu) \le n, \\ \mu \qquad \mbox{otherwise}, \end{cases}$$ 
% %$$\overline{X_{n, i}} = X_{i} 1_{d(X_{i}, \mu) \le n}:= \begin{cases} X_i \quad \mbox{if} ~ d(X_{i}, \mu) \le n, \\ \mu \qquad \mbox{otherwise}, \end{cases}$$ 
% for any $n\ge 1$ and $i=1, 2, \ldots, n$. 
%
%{\color{red}
To state our conditions,  we define, for \textit{any} reference point $\mu_0 \in M$, a triangular array of random variables $\{\overline{X_{n, i}}(\mu_0)\} _{n\ge i\ge 1}$ truncated at distance $n$ from $\mu_0$: 
$$\overline{X_{n, i}}(\mu_0) := \begin{cases} X_i \quad \mbox{if} ~ d(X_{i}, \mu_0) \le n, \\ \mu_0 \qquad \mbox{otherwise}, \end{cases}$$ 
%$$\overline{X_{n, i}} = X_{i} 1_{d(X_{i}, \mu) \le n}:= \begin{cases} X_i \quad \mbox{if} ~ d(X_{i}, \mu) \le n, \\ \mu \qquad \mbox{otherwise}, \end{cases}$$ 
for any $n\ge 1$ and $i=1, 2, \ldots, n$. 
%}

% \begin{theorem}\label{thm:main}  
% Suppose that either of the following holds:
% \begin{enumerate}
%   \item[(i)]  Condition (C1) holds, $\sum_{i=1}^{n} \mathbb{P}(d(X_i, \mu) > n) \underset{n \to \infty}{\to} 0$ and $n^{-2}\sum_{i=1}^{n}\mathbb{E}[d^2(\overline{X_{n,i}}, \mu)] \underset{n \to \infty}{\to} 0$;

%   \item[(ii)] Condition (C2) holds, and  $n\mathbb{P}(d(X_1, \mu) > n) \underset{n \to \infty}{\to} 0$.
% \end{enumerate}
% Then, $\mu_{n} \overset{\mathbb{P}}{\to} \mu$ as $n$ tends to infinity.
% \end{theorem}
\begin{theorem}\label{thm:main}  
Suppose that Condition (C) holds. If either of the following conditions is satisfied, then $\mu_{n} \overset{\mathbb{P}}{\to} \mu$ as $n$ tends to infinity.
\begin{enumerate}
  \item[(i)] For some $\mu_0 \in M$, it holds that $\sum_{i=1}^{n} \mathbb{P}(d(X_i, \mu_0) > n) \underset{n \to \infty}{\to} 0$ and $n^{-2}\sum_{i=1}^{n}\mathbb{E}[d^2\{\overline{X_{n,i}}(\mu_0), \mu_0\}] \underset{n \to \infty}{\to} 0$;

  \item[(ii)] The $X_i$'s are identically distributed and for some $\mu_0 \in M$, $n\mathbb{P}(d(X_1, \mu_0) > n) \underset{n \to \infty}{\to} 0$.
\end{enumerate}

\end{theorem}





Theorem~\ref{thm:main}  generalizes the sufficient direction of the classical result  (\ref{KF_wlln}) to random variables on non-compact symmetric spaces. 
However, their forms appear somewhat different. This discrepancy arises because while Theorem~\ref{thm:main} is stated under a symmetry assumption on $\mathbb{P}_{X_i}$, (\ref{KF_wlln}) remains valid even when $\mathbb{P}_{X}$ is non-symmetric. Notably, under the condition of (\ref{KF_wlln}), there are examples where $\mathbb{P}_{X}$ is not symmetric and $\mathbb{E}[X_1 1_{|X_1|\le n}]$ diverges as $n\to \infty$. Therefore, the classical case (\ref{KF_wlln}) includes the case where the sample mean $\bar{X}_n$ diverges. %Such a diverging sample mean is prevented when $\mathbb{P}_{X}$ is symmetric. In particular, 
In contrast, if $\mathbb{P}_{X}$ defined on $\mathbb{R}$ is symmetric about $\mu$ as we have assumed in Theorem~\ref{thm:main}, then $\mathbb{E}[X_1 1_{|X_1|\le n}] \underset{n\to \infty}{\to} \mu$. 


 
While our proof of Theorem \ref{thm:main} is provided in Appendix B, we point out three key steps in the proof. First, we observe and use the fact that if the tail conditions such as $\lim_{n\to\infty} n \mathbb{P}(d(X_1,\mu_0)>n)=0$ are satisfied for some $\mu_0\in M$, then they are satisfied for the center of symmetry $\mu$. Second, in the evaluation of $\mathbb{P}(d(\mu_n,\mu) > \epsilon)$, for a given $\epsilon>0$, the sample \Frechet\ mean $\mu_n$ is replaced by the sample \Frechet\ mean $\bar\mu_n = \mu_n\big(\overline{X_{n, 1}}(\mu), \overline{X_{n, 2}}(\mu), \ldots, \overline{X_{n, n}}(\mu)\big)$ of truncated random variables. Finally, we establish and utilize the fact that on a Hadamard space, the sample \Frechet\ mean $\bar\mu_n$ tends to shrink towards the center $\mu$ of symmetry, compared to a Euclidean counterpart. In fact, we establish this contraction argument for any deterministic data set $\{x_1, x_2, \ldots,x_n\} \subset M$ and for \textit{any base point} $x \in M$, which may not be the same as $\mu$. 

%A key technical ingredient in the proof of Theorem \ref{thm:main} is the following observation. For a base point $x$ lying on $M$ that is Hadamard, we obtain an upper bound on the distance between the sample \Frechet\ mean and the base point. Specifically, the next lemma shows that the sample \Frechet\ mean on Hadamard spaces tends to shrink towards the base point, compared to Euclidean cases.

\begin{lemma}\label{lem:Hada}
Suppose that $(M,d)$ is a Hadamard manifold. For any $x,x_1, x_2,\ldots,x_n \in M$, 
\begin{equation}\label{ineq:key}
\|\mbox{Log}_{x}(\mu_{n}) \|_{x} \le  \|\frac{1}{n}\sum_{i=1}^{n}\mbox{Log}_{x}(x_i) \|_{x},
\end{equation}
where $\mu_{n}$ denotes the sample \Frechet\ mean of $(x_1, x_2, \ldots x_n) \in M^n$, and $\|\cdot\|_{x}$ denotes the Riemannian norm induced by the Riemannian metric $\langle ~,~\rangle_{x}$ on $T_{x}M$. The inequality in (\ref{ineq:key}) is strict unless $x,x_1,…,x_n$ all lie in a single totally geodesic flat submanifold of $M$.
%unless either $M$ is flat (curvature zero everywhere) or $x, x_1, x_2, \ldots x_n$ all lie in a single geodesic. 
\end{lemma}

%The contraction property in Lemma~\ref{lem:Hada} is not an immediate consequence of \citep[][Theorem 6.3]{sturm2003probability}. Although both results embody a contraction phenomenon in non-positively curved spaces, the present inequality (\ref{ineq:key}) is expressed through the linear structure of the tangent space via the logarithmic map, whereas Sturm's theorem is stated entirely in metric terms.


% We next give a proof of Lemma~\ref{lem:Hada}. 

% \begin{proof}[of Lemma~\ref{lem:Hada}]
% If $x=\mu_{n}$, equation (2) in the main paper is true since its left-hand side equals zero. Hence, we suppose that $x\neq \mu_n$. For a fixed $y \in M$ and for any geodesic $\gamma$ with same speed, define $f(t)=d^2(\gamma(t), y)$. Then, it is known that $f''(t) \ge 2\|\dot{\gamma}(t)\|^2_{\gamma(t)}$ when $M$ is a Hadamard manifold (see, e.g., \cite{do1992riemannian}). This inequality becomes strict when $M$ has negative curvatures everywhere and $y$ does not lie in $\gamma$. To employ the inequality, for each $i=1, 2, \ldots, n$, define $g_{i}(t) = d^2(\gamma(t), x_i)$ where $\gamma$ denotes the geodesic joining $x ~ (=\gamma(0))$ and $\mu_n ~(=\gamma(1))$ with same speed. Since $\gamma$ has the same speed, $\|\dot{\gamma}(t)\|_{\gamma(t)}=\|\mbox{Log}_{x}\mu_{n}\|_{x}$ for any $0\le t \le 1$. An application of mean-value theorem leads that for each $1\le i\le n$, $g'_i(1)-g'_i(0) \ge 2\|\mbox{Log}_{x}\mu_{n}\|^2_{x}$. The summation of the previous inequality from $i=1$ to $n$ gives
% \begin{equation}\label{ineq:hadaineq}
% \sum_{i=1}^{n} \{g'_{i}(1) - g'_{i}(0)\} \ge 2n\|\mbox{Log}_{x}(\mu_n) \|^2_{x}.
% \end{equation}
% For each $1\le i\le n$, differentiating each function $g_i(t)$, we obtain
% \begin{align*}
% g'_i(0)&= -2\langle\mbox{Log}_{x}(x_i),~ \dot{\gamma}(0) \rangle_{x}= -2\langle\mbox{Log}_{x}(x_i),~ \mbox{Log}_{x}(\mu_n)\rangle_{x}~ \mbox{and} \\
% g'_{i}(1) &= -2\langle\mbox{Log}_{\mu_n}(x_i),~ \dot{\gamma}(1)\rangle_{\mu_n} = 2\langle\mbox{Log}_{\mu_n}(x_i),~ \mbox{Log}_{\mu_n}(x)\rangle_{\mu_n}. 
% \end{align*}
% In particular, we get $\sum_{i=1}^{n}g'_{i}(1)=\langle 2\sum_{i=1}^{n}\mbox{Log}_{\mu_n}(x_i),~ \mbox{Log}_{\mu_n}(x)\rangle_{\mu_n} = 0$, since $$\textsf{grad}|_{z=\mu_n}\{\sum_{i=1}^{n} d^2(z, x_i) \} = -2 \sum_{i=1}^{n} \mbox{Log}_{\mu_n}(x_i)=\textbf{0}.$$ With this in mind, according to (\ref{ineq:hadaineq}),
% \begin{equation}\label{eq:log_inner}
% -\sum_{i=1}^{n} g'_i(0)=2 \langle\sum_{i=1}^{n}\mbox{Log}_{x}(x_i),~ \mbox{Log}_{x}(\mu_{n}) \rangle_{x} \ge 2n\|\mbox{Log}_{x}(\mu_n) \|^2_{x}.
% \end{equation}
% Applying the Cauchy-Schwarz inequality, we obtain
% \begin{equation}\label{eq:cs_ineq}
% \|\sum_{i=1}^{n}\mbox{Log}_{x}(x_i) \|_{x} \cdot \|\mbox{Log}_{x}(\mu_n) \|_{x} \ge \langle \sum_{i=1}^{n}\mbox{Log}_{x}(x_i),~ \mbox{Log}_{x}(\mu_{n}) \rangle_{x}.
% \end{equation}
% Combining (\ref{eq:log_inner}) and (\ref{eq:cs_ineq}), we complete the proof. 

% Note that equality in (\ref{eq:cs_ineq}) holds when $x,x_k,\mu_n$ are positioned on the same geodesic. In this regard, the inequality in equation (2) becomes strict when $M$ is a manifold with negative curvature and $x, x_{1}, x_{2}, \ldots, x_{n}$ are not located on a common geodesic.
% \end{proof}

 % \begin{remark}
 % Lemma \ref{lem:Hada} is of independent interest. In particular, it can be used to show that the rate of convergence of the sample \Frechet\ mean is faster than in Euclidean spaces. See Remark A1 in the supplementary material for details and for its connection to variance modulation \citep{pennec2019curvature, hundrieser2024finite}.
 % \end{remark}

An important question related to Theorem~\ref{thm:main} is whether the
converse of case~(ii) holds. This concerns the extent to which the Kolmogorov--Feller tail condition is tight for the weak law. Under Condition~(C), the converse holds in the flat case, that is, for non-compact symmetric spaces with zero curvature; see Appendix D. In general, however, we conjecture that the converse need not hold. A converse would follow if the inequality (\ref{ineq:key}) could be reversed, up to a multiplicative constant. Such a reversal may not be available in negatively curved spaces, where curvature effects may substantially alter the behavior of \Frechet\ means. We therefore leave the problem of proving a general nonlinear converse theorem, or constructing a counterexample, for future work.
Another potential direction for future work is to investigate Kolmogorov--Feller weak laws of large numbers for dependent random variables in the present non-compact symmetric setting, in
analogy with \cite{naderi2020weak}. 

%Before closing the paper, we remark the following. As future work, one may investigate the Kolmogorov--Feller weak law of large numbers for dependent random variables in the present non-compact symmetric setting, in analogy with \cite{naderi2020weak}.

\section{Examples}

We present examples where our result applies to i.i.d. sequences of random variables. 

\begin{example}[Symmetric positive-definite matrices]
Let \(M=\mathrm{Sym}^+(k)\), the space of $k \times k$ symmetric positive-definite matrices, be equipped with the affine-invariant Riemannian distance
\[
d_{\mbox{\tiny AI}}(A,B)=\|\log(A^{-1/2}BA^{-1/2})\|_{F}, \qquad (A,B \in M)
\]
where $\|\cdot\|_{F}$ stands for the Frobenius norm. Then $(M, d_{\mbox{\tiny AI}})$ is a non-compact symmetric space with non-positive curvature. At \(\mu=I_{k}\), the $k \times k$ identity matrix, the geodesic symmetry is given by $s_{\mu}(A)=A^{-1}$. Hence, a $k\times k$ positive-definite random matrix $X$ satisfying $X\stackrel{d}{=}X^{-1}$ is geodesically symmetric about $\mu$. In this setting, let $X, X_1, X_2, \ldots$ be a sequence of i.i.d. random variables. Theorem~\ref{thm:main}(ii) gives that for any $\mu_0 \in M$, $n\mathbb{P}(d_{\mbox{\tiny AI}}(X, \mu_0) > n) \underset{n \to \infty}{\to} 0$ implies $\mu_n(X_1, \ldots, X_n) \overset{\mathbb{P}}{\to} \mu$ as $n\to\infty$. 
\end{example}

\begin{example}[Product symmetric spaces]
Let $(M_j, d_j)_{j=1}^{K}$ be non-compact symmetric spaces and equip $M=M_1\times\cdots\times M_K$ with the product distance $$d(x,y)=\{\sum_{j=1}^K d_j^2(x_j,y_j)\}^{1/2}$$ for any $x=(x_1, \ldots, x_K)$ and $y=(y_1, \ldots, y_K)$. Then, $(M,d)$ is again a non-compact symmetric space. For a $\mu=(\mu_1,\ldots,\mu_K) \in M$, the geodesic symmetry about \(\mu\) is given by
\[
s_\mu(x_1,\ldots,x_K)
=
(s_{\mu_1}(x_1),\ldots,s_{\mu_{K}}(x_K)),
\]
where $s_{\mu_j}$ denotes the geodesic symmetry about $\mu_j \in M_j$. Let $\pi_j:M \to M_j, ~ 1\le j\le K$, be the natural projection map. Theorem~\ref{thm:main}(ii) applies to product-valued data whenever $X\stackrel{d}{=}s_\mu(X)$ and
$n\mathbb P(d(X,\mu)>n)\underset{n\to \infty}{\to} 0.$ The product-distance tail condition is equivalent to the coordinate-wise tail conditions $n\mathbb P\{d_{j}(\pi_{j}(X),\mu_j)>n\} \underset{n\to \infty}{\to} 0, \quad j=1, \ldots, K.$ Marginal geodesic symmetry alone does not imply joint geodesic symmetry in general, since the dependence structure may break the symmetry. However, if $\pi_1(X), \ldots, \pi_K(X)$ are independent, then $X\stackrel{d}{=}s_\mu(X)$ is equivalent to $\pi_j(X)\stackrel{d}{=}s_{\mu_{j}}(\pi_j(X))$ for all $1\le j\le K$. Thus, under independence, the geodesic symmetry condition can be verified coordinate-wise.  
\end{example}




%%
We next present an illustrative example where our result applies to a non-i.i.d. sequence of random variables $(X_i)_{i=1}^{\infty}$ in $\mathrm{Sym}^+(k)$.  
%Specifically, we consider independent random variables $(X_i)_{i=1}^{\infty}$ taking values in $\mathrm{Sym}^+(k)$.
In this example, although each $X_i$ has finite \Frechet\ variance $\sigma_i^2 := \inf_{p\in M} \mathbb{E}[d^2(X_i, p)]$, the sequence of variances $(\sigma_i^2)_{i=1}^\infty$ is unbounded. 
%} 

%We note that the hypothesis of Theorem~\ref{thm:main} for the case (i) holds when $(X_i)_{i=1}^{\infty}$ does not necessarily consist of identically distributed random variables. An illustrative example with non-i.i.d. random variables in the space of symmetric positive-definite matrices is provided below.

\begin{example} 
Consider a sequence $(X_i)_{i=1}^{\infty}$ of independent random variables taking values in $(\mathrm{Sym}^+(k),d_{\text{\tiny AI}})$ for some $k \ge 1$. Let $\mu = I_k$, the $k \times k$ identity matrix, be the center of symmetry. For a fixed $\alpha \in (0,1)$, each $X_i$ follows a log-Gaussian distribution:
$$\textsf{vec}\{\mbox{Log}_{\mu}(X_i)\} \sim N_{k(k+1)/2}(\textbf{0}, ~i^{\alpha}I_{k(k+1)/2}),$$ 
where $\textsf{vec}(\cdot)$ denotes the vectorization operator appropriate for symmetric $k \times k$ matrices \citep[][Section 3.5]{pennec2006riemannian}.  
The log-Gaussian distributions are geodesically symmetric about $\mu$, which is also the  \Frechet\ mean of each $X_i$, while the \Frechet\ variance $\sigma_i^2 = i^\alpha k(k+1)/2$ diverges as $i$ increases. 
In this setting, the conditions in Theorem~\ref{thm:main}(i) are satisfied (see Appendix C for details), and we conclude that $\mu_n(X_1,\ldots, X_n) \to \mu$ in probability as $n \to \infty$. 
%Since the conditions of Theorem~\ref{thm:main}(i) are satisfied, we can conclude that $\mu_n(X_1, X_2, \ldots , X_n) \overset{\mathbb{P}}{\to} \mu$ as $n$ tends to infinity.
\end{example}



% As a corollary of Theorem~\ref{thm:main}, the weak law of large numbers for i.i.d. sequences with finite first moment, $\mathbb{E}[d(X_1, \mu)]<\infty$, follows immediately.  
% \begin{corollary}[Weak law of large numbers]\label{cor:wlln} 
% Suppose that Condition (C2) holds and the first moment of $X_1$ exists; that is, $\mathbb{E}[d(X_1, \mu)] < \infty$. Then, $\mu_n \overset{\mathbb{P}}{\to} \mu$ as $n$ goes to infinity.
% \end{corollary}

% \begin{proof}[Proof of Corollary \ref{cor:wlln}]
% According to Theorem~\ref{thm:main} for the case (ii), it suffices to show that $n\mathbb{P}(d(X_1, \mu) > n) \underset{n\to \infty}{\to} 0$. An application of the dominated convergence theorem together with $\mathbb{E}[d(X_1, \mu)] < \infty$ leads that
% $$
% n\mathbb{P}(d(X_1, \mu) > n) \le \mathbb{E}[d(X_1, \mu)\, 1_{d(X_1, \mu) > n}] \underset{n\to \infty}{\to} 0,
% $$
% yielding the desired result.
% \end{proof}

%The law of large numbers in Corollary~\ref{cor:wlln} holds for the case where the population \Frechet\ mean does not exist. In fact, such a result has been established without assuming symmetry. 
For the case where only the first moment of $\mathbb{P}_{X}$ exists, the sample \Frechet\ mean almost surely converges to the minimizer of the \Frechet\ difference $F_{\mu}(m) = \mathbb{E}[d^2(X,m) - d^2(X,\mu)]$, which is well-defined for any $m \in M$ provided that $\mathbb{E}[d(X,\mu)]<\infty$ \citep{sturm2003probability,schotz2022strong}. 
%[\cite{sturm2003probability}, Proposition 4.3., noted already that \Frechet\ mean is the minimizer of the difference as long as the first moment exists.] [While the \Frechet\ mean is defined as the minimizer of the difference, its consistency needs at least the first moment condition, for example, Theorem 3.2 and Corollary 4.3 in \cite{schotz2022strong}.] 
We emphasize that Theorem~\ref{thm:main} relaxes this finite first moment condition. 
In the absence of moment assumptions, a symmetry condition serves to identify a center of the distribution. 
%
%However, without any finite moment, the point to which the sample \Frechet\ mean converges is not well-defined. Out result does not require finite moments, but the geodesic symmetry about $\mu$ provides the point to which the sample \Frechet\ mean converges. %Our remedy is to assume geodesic symmetry of the distribution, as in Condition (C) {\color{red}and Corollary 3.3 of \cite{marcus1979stable}.} 
%under a broad class of distributions.
%%
The following example demonstrates that
the tail condition in Theorem~\ref{thm:main} is genuinely weaker than the finite
first-moment condition.
%To this end, we present an example of a distribution with infinite first moment, $\mathbb{E}[d(X, \mu)] = \infty$, that nevertheless satisfies the tail condition $\lim_{n\to\infty}n\mathbb{P}(d(X, \mu) > n) = 0$.



%{\color{blue} [\cite{schotz2022strong} assumes only $E[d(X,p)] < \infty$. We should reflect this fact. ]}
%{\color{red}We again emphasize that all existing laws of large numbers for the sample \Frechet\ mean on non-Euclidean spaces (e.g., \cite{ziezold1977expected, bhattacharya2003large, schotz2022strong, evans2024limit}) typically require the existence of second moment (or first moment for the case \cite{schotz2022strong}), i.e., $\mathbb{E}[d^2(X, m)]<\infty$ (or $\mathbb{E}[d(X, m)]<\infty$, respectively) for some (thus all) $m\in M$.
%[\cite{sturm2003probability}, Proposition 4.3., noted already that \Frechet\ mean is the minimizer of the difference as long as the first moment exists.] 
%}{\color{blue} [While the \Frechet\ mean is defined as the minimizer of the difference, its consistency needs at least the first moment condition, for example, Theorem 3.2 and Corollary 4.3 in \cite{schotz2022strong}.]} 
%
%Our results relax this requirement under a broad class of distributions.
%%
%Furthermore, we construct an example of a distribution with infinite first moment, $\mathbb{E}[d(X, \mu)] = \infty$, that nevertheless satisfies the tail condition $\lim_{n\to\infty}n\mathbb{P}(d(X, \mu) > n) = 0$, as shown below. %}
%Notably, there is at least one probability distribution $\mathbb{P}_{X}$ that does not have a finite first moment (i.e., $\mathbb{E}[d(X, \mu)] = \infty$), but satisfies $n\mathbb{P}(d(X, \mu) > n) \to 0$ as $n$ goes to infinity as follows.

\begin{example}
Suppose that $\mathbb{P}_{X}$ is a geodesically symmetric distribution about $\mu$ satisfying $$\mathbb{P}(d(X, \mu)>t) \propto \frac{1}{t\log(t)}$$ for large enough $t ~(\ge N)$. Then, it holds that $\lim_{n\to \infty}n\mathbb{P}(d(X, \mu)>n) = \lim_{n \to \infty}\frac{1}{\log(n)}= 0$. On the other hand, it is well-known that $\mathbb{E}[d(X, \mu)]=\infty$ is equivalent to $\sum_{n=N}^{\infty} \mathbb{P}(d(X, \mu) > n)= \sum_{n=N}^{\infty}\frac{c}{n\log(n)} =\infty$ for a fixed $c>0$, which follows from the integral test:
\begin{eqnarray*}
\sum_{n=N}^{\infty}\frac{1}{n\log(n)}=\infty ~ \Leftrightarrow ~ \int_{N}^{\infty}\frac{1}{x\log(x)}dx=\big[\log\{\log(x)\}\big]_{x=N}^{\infty} = \infty.
\end{eqnarray*}
Thus, by Theorem~\ref{thm:main}(ii), while the first moment does not exist for $\mathbb{P}_X$, its sample \Frechet\ mean $\mu_n$ converges to the center $\mu$ of symmetry in probability. 
%Therefore, there exists at least one probability distribution on $M$ satisfying the assumption in Theorem~\ref{thm:main} for the case (ii), but $\mathbb{E}[d(X, \mu)] = \infty$.
\end{example}



\section*{Acknowledgement}  
We are grateful to the reviewers for their thoughtful and constructive comments. Their feedback led to improvements in the exposition, examples, and positioning of the results. Jongmin Lee was supported by the National Research Foundation of Korea (NRF) grants funded by the Korea government (MSIT) (RS-2026-25479875). Sungkyu Jung was supported by the National Research Foundation of Korea (NRF) grants funded by the Korea government (MSIT) (RS-2023-00301976 and RS-2024-00333399).



% \noindent {\it Conflict of interest}: The authors do not declare any conflict of interest regarding this work.





% \section*{Funding}
% Sungkyu Jung was supported by the National Research Foundation of Korea (NRF) grants funded by the Korea government (MSIT) (RS-2023-00301976 and RS-2024-00333399). 



% \section*{Data availability}
% The multivariate tensor-based morphometry data used in this study are not publicly available. However, researchers may request access by contacting the authors of the original study \citep{paquette2017ventricular}, subject to their approval. 




\begin{appendix}

\section{Proof of Proposition 1}
\begin{proof}[Proof of Proposition 1]
Suppose, for contradiction, that $\mu_1 \neq \mu_2$. Let $T = s_{\mu_2} \circ s_{\mu_1}$, and $\gamma$ be the geodesic from $\mu_1$ to $\mu_2$. Denote $\ell = 2d(\mu_1, \mu_2) > 0$. Then, $T$ is a translation-type isometry along $\gamma$, more precisely, a transvection. (For an illustration of the transvection, see Figure~\ref{fig:transvection}.) An application of Corollary 2.5 of \cite{sturm2003probability} implies that, since $M$ is Hadamard, we have for any $x \in M$, $d(x, s_{\mu_2}\circ s_{\mu_1}(x))\ge 2d(\mu_1, \mu_2)=\ell$.
%
%It can be shown that for any $x \in M$, $d(x, s_{\mu_2}\circ s_{\mu_1}(x))\ge 2d(\mu_1, \mu_2)=\ell$ by an application of Corollary 2.5 of \cite{sturm2003probability}. 
%
This implies that, for any $x\in M$ and for any $m\in \mathbb{N}$, $d(x, T^m(x))\ge \ell m$. Thus, the isometry $T$, when repeatedly applied, pushes any point to infinity. 
%%
\begin{figure}[htbp!]
\centering
\includegraphics[scale=0.45]{figures/fig_transvection.png}
\caption{An illustration of the transvection $T$, mapping from $x$ to $s_{\mu_2}\circ s_{\mu_1}(x)$.}
\label{fig:transvection}
\end{figure}



%% 
Given bounded sets $K$ and $A$, define a constant $C_{K, A}=\sup_{k\in K, a\in A}d(k, a) < \infty$. There exists an integer $m$ such that $\ell m> C_{K,A}$. For any $k \in K$ and for any $a \in A$, by the triangle inequality, 
$$
d(k, T^{m}(a)) \ge d(a, T^{m}(a)) - d(k,a) \ge \ell m - C_{K,A} >0,
$$
which implies that $K \cap T^{m}(A) =\emptyset$, where $T^{m}(A):=\{T^m(a): a \in A\}$. That is, for any bounded sets $K$ and $A$, there exists an integer $m_0$ such that for any integer $m'\ge m_0$, $K \cap T^{m'}(A) = \emptyset$, where for any $m\ge 1$, $T^{m}(A) \subset M$.

 



%%
We now set the bounded sets $A$ and $K$ as follows. 
First, $A$ is chosen to be any bounded Borel set in $M$. Since $\mathbb{P}_X$ is geodesically symmetric about both $\mu_1$ and $\mu_2$, it holds that $\mathbb{P}_X(A) = \mathbb{P}_X(T(A))$.

Since $(M,d)$ is second countable, every probability measure on $M$ is tight \citep{lee2013smooth}. Hence, for any $\epsilon > 0$, there exists a compact set $K = K(\epsilon)$ satisfying $\mathbb{P}_X(K) \ge 1-\epsilon$. For all large $m$, we have $T^m(A) \cap K = \emptyset$, so that $\mathbb{P}(A) = \mathbb{P}(T^m(A)) \le 1 - \mathbb{P}(K) \le \epsilon$. 
%
Since $\epsilon$ is arbitrary, we have $\mathbb{P}(A) = 0$ for any bounded Borel set $A \subset M$. This is a contradiction to the tightness of $\mathbb{P}_{X}$, and therefore, $\mu_1 = \mu_2$.
\end{proof}


\section{Proofs of Theorem 1 and Lemma 1, and auxiliary results}
Our proof of Theorem 1 relies on several lemmas. We first state a result on the tail conditions imposed in Theorem 1. 

\begin{lemma}\label{lem:somevsall_tail_conditions} 
Let $(M,d)$ be a metric space, and let $(X_i)_{i=1}^\infty$ be a sequence of independent $M$-valued random variables. For an $m \in M$, let $c_1(m)$, $c_2(m)$ and $c_3(m)$ denote the following conditions, respectively: 
\begin{align*}
   c_1(m): & \lim_{n\to\infty} n\mathbb{P}(d(X_1,m) > n) = 0;\\
   c_2(m): & \lim_{n\to\infty} n^{-2}\sum_{i=1}^n\mathbb{E}[d^2\{\overline{X_{n,i}}(m),m\}] = 0;\\
   c_3(m): & \lim_{n\to\infty} \sum_{i=1}^n\mathbb{P}(d(X_i,m) > n) = 0.
\end{align*}
\begin{itemize}
    \item[(i)] For any $m,m' \in M$, $c_1(m)$ holds if and only if $c_1(m')$ holds.
    \item[(ii)] For any $m,m' \in M$, $c_2(m)$ holds if and only if $c_2(m')$ holds.
    \item[(iii)] Additionally assume that $M$ is a symmetric space and,  for a given $\mu \in M$, $\mathbb{P}_{X_i}$ is geodesically symmetric about $\mu$. Then, for any $m \in M$, $c_3(\mu)$ holds if $c_3(m)$ holds.
\end{itemize}
\end{lemma}


\begin{proof}[Proof of Lemma \ref{lem:somevsall_tail_conditions}]
(i) Let $m, m'$ be fixed, and without losing generality, assume that $R:= d(m,m') > 0$. We shall show that the condition $c_1(m)$ implies $c_1(m')$. The converse is obtained by the symmetry of the argument. 

Let $x \in M$ be any point. 
By the triangle inequality, 
\begin{equation}
   d(x,m') \le d(x, m) + d(m,m') = d(x, m) + R. \label{eq:triangle}
\end{equation} 
Therefore, for any $n > R$, $d(x,m') > n$ implies $d(x,m)> n - R$, which in turn implies that 
$$\{x \in M: d(x,m') > n\} \subset \{x \in M: d(x,m)> n - R\}.$$
Then, 
\begin{align*}
    n\mathbb{P}(d(X_1,m')>n) & \le n\mathbb{P}(d(X_1,m)>n-R) \\
      & = \frac{n}{n-R} (n-R) \mathbb{P}(d(X_1,m)>n-R).
\end{align*}
But, since we assumed $c_1(m)$, we have $(n-R) \mathbb{P}(d(X_1,m)>n-R) \to 0$ as $n\to\infty$, and $\lim_{n\to\infty}{n}/({n-R})= 1$ as $R$ is fixed. Thus, $n\mathbb{P}(d(X_1,m')>n) \to 0$ as $n\to\infty$, which is precisely the condition $c_1(m')$. 

(ii) Let $m, m'$ be fixed, and let $R:= \lceil d(m,m') \rceil > 0$ be the smallest non-zero integer larger than or equal to $d(m,m')$. As in Part (i), it suffices to show that $c_2(m)$ implies $c_2(m')$. 
Observe first that, for any $X_i,m \in M$, 
$$d^2(\overline{X_{n,i}}(m),m) = d^2(X_i,m) 1_{d(X_i,m) \le n},$$
which can be deduced from the definition of $\overline{X_{n,i}}(m)$. 

By the triangle inequality %(\ref{eq:triangle})
we have 
$$\{x \in M: d(x,m') \le n\} \subset \{x \in M: d(x,m)\le n + R\},$$ 
and, for any $x \in M$ 
$$d^2(x,m') \le \{d(x,m) + R\}^2 \le 2d^2(x,m) + 2R^2.$$
It follows that
\begin{align*}
    d^2(X_i,m') 1_{d(X_i,m') \le n} 
    & \le [2d^2(X_i,m) + 2R^2]1_{d(X_i,m') \le n} \\ 
    & \le [2d^2(X_i,m) + 2R^2]1_{d(X_i,m)\le n + R}.
\end{align*}

Write $F_m(n) : = n^{-2}\sum_{i=1}^n\mathbb{E}[d^2\{\overline{X_{n,i}}(m),m\}]$ for any $m\in M$ and $n\ge 1$. Then, 
\begin{align*}
    F_{m'}(n) & = n^{-2}\sum_{i=1}^n\mathbb{E} d^2(X_i,m') 1_{d(X_i,m') \le n}  \\
    & \le \frac{2}{n^2}\sum_{i=1}^n\mathbb{E} d^2(X_i,m)1_{d(X_i,m)\le n + R} + \frac{2R^2}{n^2}\sum_{i=1}^n \mathbb{P}(d(X_i,m)\le n + R) \\ 
    & \le \frac{2(n+R)^2}{n^2}F_{m}(n+R) + \frac{2R^2}{n^2}\sum_{i=1}^n 1.
\end{align*}
Since $F_{m}(n+R) \to 0$ as $n\to \infty$ (which is precisely the condition $c_2(m)$), $F_{m'}(n) \underset{n\to \infty}{\to} 0$ as desired. 





(iii) Recall that $s_\mu(\cdot)$ is the geodesic symmetry with respect to $\mu \in M$. For any $x \in M$, we have 
$d(x, s_\mu(x)) = 2d(x,\mu) = 2d(s_\mu(x),\mu)$ (see Figure~\ref{fig:transvection} for an illustration). By the triangle inequality, for any $m\in M$, we obtain 
\begin{equation}
    \label{eq:iii} 
    2d(x,\mu) = d(x, s_\mu(x)) \le d(x,m ) + d(m, s_\mu(x)).
\end{equation}

Let us now fix $m$, and suppose that $c_3(m)$ holds. 
From (\ref{eq:iii}), we have 
\begin{align*}
\{x\in M: d(x,\mu) > n\} &\subset \{x\in M: d(x,m) + d(s_\mu(x), m) > 2n\}\\
&\subset \{x\in M: d(x,m) > n\}\cup \{x\in M: d(s_\mu(x),m) > n\}.
\end{align*}
Since $\mathbb{P}_{X_i}$ is geodesically symmetric about $\mu$, $X_i$ and $s_\mu (X_i)$ have the same distribution. Therefore, we have 
$$
\mathbb{P}(d(X_i,\mu) > n ) \le 2\mathbb{P}(d(X_i,m) > n).$$
Summing over $1\le i\le n$, we have 
$$
\sum_{i=1}^n\mathbb{P}(d(X_i,\mu) > n ) \le 2\sum_{i=1}^n\mathbb{P}(d(X_i,m) > n) \to 0,
$$
as $n\to \infty$. That is, $c_3(\mu)$ holds.
\end{proof}


%\begin{remark}
The first two results of Lemma~\ref{lem:somevsall_tail_conditions} are valid for any metric space, and no symmetry condition is required. In contrast, the geodesic symmetry condition is required for Lemma~\ref{lem:somevsall_tail_conditions}(iii), and the converse is generally not true. That is, while $c_3(\mu)$ is true whenever $c_3(m)$ is true for some $m$, $c_3(m')$ may not be true even if $c_3(\mu)$ is true. A simple counterexample is the following: let $(M,d)$ be the usual real line with standard metric, and let $X_i = i$ with probability $1/2$ and $-i$ with probability $1/2$. The center of symmetry for $\mathbb{P}_{X_i}$ is $\mu = 0$. It can be checked that $\sum_{i=1}^n \mathbb{P}(d(X_i,\mu) > n) = \sum_{i=1}^n \mathbb{P}(|X_i - \mu| > n)= 0$ for all $n\ge 1$, but for $m' = 1$, $\sum_{i=1}^n \mathbb{P}(d(X_i,m') > n) = 1/2$ for all $n\ge 1$. 
%\end{remark}

Thanks to Lemma \ref{lem:somevsall_tail_conditions}, we shall show Theorem 1 under the following simplified condition:
\begin{enumerate}
  \item[(i)]  Condition (C) holds, $\sum_{i=1}^{n} \mathbb{P}(d(X_i, \mu) > n) \underset{n \to \infty}{\to} 0$ and \linebreak $n^{-2}\sum_{i=1}^{n}\mathbb{E}[d^2(\overline{X_{n,i}}, \mu)] \underset{n \to \infty}{\to} 0$;

  \item[(ii)] Condition (C) holds, the $X_i$'s are identically distributed, and it holds that $n\mathbb{P}(d(X_1, \mu) > n) \underset{n \to \infty}{\to} 0$.
\end{enumerate}

 

We next give a proof of Lemma 1. 

\begin{proof}[Proof of Lemma 1]
If $x=\mu_{n}$, equation (2) in the main paper is true since its left-hand side equals zero. Hence, we suppose that $x\neq \mu_n$. For a fixed $y \in M$ and for any geodesic $\gamma$ with same speed, define $f(t)=d^2(\gamma(t), y)$. Then, it is known that $f''(t) \ge 2\|\dot{\gamma}(t)\|^2_{\gamma(t)}$ when $M$ is a Hadamard manifold (see, e.g., \cite{do1992riemannian}). This inequality becomes strict when $M$ has negative curvatures everywhere and $y$ does not lie in $\gamma$. To employ the inequality, for each $i=1, 2, \ldots, n$, define $g_{i}(t) = d^2(\gamma(t), x_i)$ where $\gamma$ denotes the geodesic joining $x ~ (=\gamma(0))$ and $\mu_n ~(=\gamma(1))$ with same speed. Since $\gamma$ has the same speed, $\|\dot{\gamma}(t)\|_{\gamma(t)}=\|\mbox{Log}_{x}\mu_{n}\|_{x}$ for any $0\le t \le 1$. An application of mean-value theorem implies that for each $1\le i\le n$, $g'_i(1)-g'_i(0) \ge 2\|\mbox{Log}_{x}\mu_{n}\|^2_{x}$. The summation of the previous inequality from $i=1$ to $n$ gives
\begin{equation}\label{ineq:hadaineq}
\sum_{i=1}^{n} \{g'_{i}(1) - g'_{i}(0)\} \ge 2n\|\mbox{Log}_{x}(\mu_n) \|^2_{x}.
\end{equation}
For each $1\le i\le n$, differentiating each function $g_i(t)$, we obtain
\begin{align*}
g'_i(0)&= -2\langle\mbox{Log}_{x}(x_i),~ \dot{\gamma}(0) \rangle_{x}= -2\langle\mbox{Log}_{x}(x_i),~ \mbox{Log}_{x}(\mu_n)\rangle_{x}~ \mbox{and} \\
g'_{i}(1) &= -2\langle\mbox{Log}_{\mu_n}(x_i),~ \dot{\gamma}(1)\rangle_{\mu_n} = 2\langle\mbox{Log}_{\mu_n}(x_i),~ \mbox{Log}_{\mu_n}(x)\rangle_{\mu_n}. 
\end{align*}
In particular, we get $\sum_{i=1}^{n}g'_{i}(1)=\langle 2\sum_{i=1}^{n}\mbox{Log}_{\mu_n}(x_i),~ \mbox{Log}_{\mu_n}(x)\rangle_{\mu_n} = 0$, since $$\textsf{grad}|_{z=\mu_n}\{\sum_{i=1}^{n} d^2(z, x_i) \} = -2 \sum_{i=1}^{n} \mbox{Log}_{\mu_n}(x_i)=\textbf{0}.$$ With this in mind, according to (\ref{ineq:hadaineq}),
\begin{equation}\label{eq:log_inner}
-\sum_{i=1}^{n} g'_i(0)=2 \langle\sum_{i=1}^{n}\mbox{Log}_{x}(x_i),~ \mbox{Log}_{x}(\mu_{n}) \rangle_{x} \ge 2n\|\mbox{Log}_{x}(\mu_n) \|^2_{x}.
\end{equation}
Applying the Cauchy-Schwarz inequality, we obtain
\begin{equation}\label{eq:cs_ineq}
\|\sum_{i=1}^{n}\mbox{Log}_{x}(x_i) \|_{x} \cdot \|\mbox{Log}_{x}(\mu_n) \|_{x} \ge \langle \sum_{i=1}^{n}\mbox{Log}_{x}(x_i),~ \mbox{Log}_{x}(\mu_{n}) \rangle_{x}.
\end{equation}
Combining (\ref{eq:log_inner}) and (\ref{eq:cs_ineq}), we complete the proof. 

Note that equality in (\ref{eq:cs_ineq}) holds when $x,x_k,\mu_n$ are positioned on the same geodesic. In this regard, the inequality in equation (2) becomes strict when $M$ is a manifold with negative curvature and $x, x_{1}, x_{2}, \ldots, x_{n}$ are not located on a common geodesic.
\end{proof}


The following lemma serves as a step toward establishing the main results under Condition (C). 
\begin{lemma}\label{lem:modulation_ineq} 
Suppose that $(M,d)$ is a non-compact symmetric space and Condition (C) holds.
\begin{enumerate}
    \item[(a)] If, for each $1\le i\le n$, $\mathbb{P}_{X_{i}}$ has finite variance, i.e., $\mathbb{E}[d^2(X_{i}, \mu)] < \infty$, then 
     \begin{equation}\label{ineq:mod}
        \mathbb{E}[d^2(\mu_n, \mu)]\le \frac{1}{n^2}\sum_{i=1}^{n}\mathbb{E}[d^2(X_i, \mu)].
     \end{equation}

    \item[(b)] In particular, if $X_i$'s are identically distributed and $\mathbb{P}_{X}$ has finite variance, then
    \begin{equation}\label{ineq:modulation}
        \mathbb{E}[d^2(\mu_n, \mu)]\le \frac{1}{n}\mathbb{E}[d^2(X_1, \mu)].
    \end{equation}
    \end{enumerate}
\end{lemma}


\begin{proof}[Proof of Lemma~\ref{lem:modulation_ineq}]
(a) Since $(M,d)$ is a non-compact symmetric space, it is simply connected (see, e.g., \cite{helgason1979differential}). Thus, $M$ is a Hadamard manifold. Hence, we can employ the result of Lemma 1, which is presented in the main paper. Notably, the population \Frechet\ mean with respect to $\mathbb{P}_{X}$ is unique and equals $\mu$ since $M$ is a non-compact symmetric space (for details, see Proposition 2(a) in \cite{lee2026huber}). Thus, for each $i=1, 2, \ldots, n$, it holds that $\mathbb{E}[\mbox{Log}_{\mu}(X_i)] = \textbf{0} \in T_{\mu}M ~(\cong\mathbb{R}^{k})$ (since each $\mathbb{P}_{X_i}$ has finite variance). Substituting $x=\mu$, and $x_i=X_i$ for any $i=1,2,\ldots n$ into equation (2), we get
$$
\|\mbox{Log}_{\mu}(\mu_n) \|_{\mu} \le \|\frac{1}{n}\sum_{i=1}^{n} \mbox{Log}_{\mu}(X_i) \|_{\mu}.
$$
Squaring both sides of the above inequality and taking the expectation, we obtain the desired result.


(b) Since $X_{1}, X_2, \ldots, X_n$ has same distribution $\mathbb{P}_{X}$ by assumption, for any $1\le i\le n$, it holds that $\mathbb{E}[d^2(\mu, X_1)]=\mathbb{E}[d^2(\mu, X_i)]$, yielding (\ref{ineq:modulation}). 


Notably, the inequality in (\ref{ineq:modulation}) becomes strict when $(M,d)$ has negative curvatures everywhere and $\mathbb{P}_X$ is not supported entirely on any single geodesic.
\end{proof}



 \begin{remark} 
         Suppose that Condition (C) holds, $X_i$ are identically distributed, and that $\mathbb{P}_{X}$ has finite variance. Then, (\ref{ineq:modulation}) implies that the variance modulation of \Frechet\ mean (see, for example, \cite{pennec2019curvature} and eq. (1.5) in \cite{hundrieser2024finite} for details), denoted by $\textsf{m}_{n} \ge 0$,
         is not greater than 1; that is, 
 \[
 \textsf{m}_{n}^2:=\frac{n\mathbb{E}[d^2(\mu_n, \mu)]}{\mathbb{E}[d^2(X, \mu)]} \le 1
 \]
 where $\textsf{m}_n < 1$ holds when $M$ has negative curvatures everywhere and $\mathbb{P}_{X}$ is not supported entirely on any single geodesic. The strict inequality $\textsf{m}_n < 1$ means that the rate of convergence for \Frechet\ mean is faster than the Euclidean counterpart. This result is more general than the existing one \citep[][Theorem 9]{pennec2019curvature} (when $M$ has constant negative curvatures), as it applies to a broader class of manifolds $M$ and underlying distributions $\mathbb{P}_{X}$ on $M$.
 \label{rmk:modulation}
 \end{remark}



\begin{lemma}\label{lem:moment}
For a real-valued random variable $Y$ and a fixed $p > 0$, 
$$
\mathbb{E}|Y|^{p}=\int_{0}^{\infty} pt^{p-1}\mathbb{P}(|Y|>t)dt.
$$
\end{lemma}
\begin{proof}[Proof of Lemma~\ref{lem:moment}]
For any non-negative random variable $Z$, it holds that $\mathbb{E}Z=\int_{0}^{\infty} \mathbb{P}(Z>t)dt$. By plugging-in $Z=|Y|^{p-1}$ into the previous equality,
\begin{align*}
\mathbb{E}|Y|^{p} &= \int_{0}^{\infty}\mathbb{P}(|Y|^p>u)du = \int_{0}^{\infty}\mathbb{P}(|Y|^p>t^p)pt^{p-1}dt \\
&=\int_{0}^{\infty} pt^{p-1}\mathbb{P}(|Y|>t)dt,
\end{align*}
where the second equality holds due to the integration by substitution for $u=t^{p}$.
\end{proof}
We are now ready to prove Theorem 1.


\begin{proof}[Proof of Theorem 1]

\textbf{Case I.} We begin by showing that if (i) in the statement of Theorem 1 holds, then $\mu_n \overset{\mathbb{P}}{\to} \mu$ as $n \to \infty$. Denote the sample \Frechet\ mean of $(\overline{X_{n, i}})_{1\le i\le n}$ by $\overline{\mu_{n}}$. Note that, for each $1\le i\le n$, the distribution of $\overline{X_{n, i}}$ is geodesically symmetric about $\mu$ and has finite variance. For a given $\epsilon > 0$, an application of Lemma~\ref{lem:modulation_ineq} implies that
\begin{eqnarray*}
\mathbb{P}(d(\mu_n, \mu ) > \epsilon) &\le& \mathbb{P}(\mu_n \neq \overline{\mu_n}) + \mathbb{P}(\mu_n=\overline{\mu_n}, ~ d(\mu_{n}, \mu) > \epsilon) \\
&\le& \mathbb{P}(\overline{X_{n, i}} \neq X_i ~ \mbox{for some} ~ i) + \mathbb{P}(d(\overline{\mu_n}, \mu) > \epsilon) \\
&\le& \sum_{i=1}^{n}\mathbb{P}(d(X_i, \mu) > n) + \mathbb{P}(d(\overline{\mu_n}, \mu) > \epsilon). \\
&=& \sum_{i=1}^{n}\mathbb{P}(d(X_i, \mu) > n) + \mathbb{E}1_{d(\overline{\mu_n}, \mu) > \epsilon} \\
&\le& \sum_{i=1}^{n}\mathbb{P}(d(X_i, \mu) > n) + \mathbb{E}\big[\frac{d(\overline{\mu_{n}}, \mu)}{\epsilon}\big]^2 \\
&\overset{(\because ~ Lemma~\ref{lem:modulation_ineq}(a))}{\le}& \sum_{i=1}^{n}\mathbb{P}(d(X_i, \mu) > n) + \frac{1}{n^2\epsilon^2}\sum_{i=1}^{n}\mathbb{E}[d^2(\overline{X_{n, i}}, \mu)] \\
&\underset{n \to \infty}{\to}& 0,
\end{eqnarray*}
where the last holds due to the hypothesis of Theorem 1 for the case (i). Therefore, $\mu_n$ converges in $\mu$ in probability as $n\to \infty$, as asserted. 


\textbf{Case II.} We next prove that if the hypothesis of Theorem 1 for the case (ii) holds, then $\mu_n \overset{\mathbb{P}}{\to} \mu$ as $n \to \infty$. By assumption, for each $n\ge 1$, $X_{1}, X_{2}, \ldots X_{n}$ have a common distribution, and $\overline{X_{n,1}}, \overline{X_{n,2}}, \ldots, \overline{X_{n,n}}$ are also identically distributed. Due to \textbf{Case I}, it is enough to show that $\sum_{i=1}^{n} \mathbb{P}(d(X_i, \mu) > n) \underset{n \to \infty}{\to} 0$ and $n^{-2}\sum_{i=1}^{n}\mathbb{E}[d^2(\overline{X_{n, i}}, ~ \mu)] \underset{n \to \infty}{\to} 0$ are satisfied. 

It directly follows that $\sum_{i=1}^{n} \mathbb{P}(d(X_i, \mu) > n) = n \mathbb{P}(d(X_1, \mu) > n)\underset{n \to \infty}{\to} 0$. An application of Lemma~\ref{lem:moment} for $p=2$ implies that for any real-valued random variable $Y$, $\mathbb{E}(Y^2)=\int_{0}^{\infty} 2t\,\mathbb{P}(|Y|>t)dt$. With this in mind, observe that
\begin{align}\label{eq:int}
n^{-2}\sum_{k=1}^{n}\mathbb{E}[d^2(\overline{X_{n,k}},\mu)] &= n^{-1}\mathbb{E}[d^2(\overline{X_{n,1}}, \mu)] \nonumber \\
&=n^{-1} \int_{0}^{\infty} 2t\, \mathbb{P}(d(\overline{X_{n,1}}, \mu) > t)dt \nonumber \\
&\le 2n^{-1}\int_{0}^{n}t\, \mathbb{P}(d(X_1, \mu) > t)dt 
\end{align}
The remainder of this proof is to verify that (\ref{eq:int}) goes to zero as $n \to \infty$. To do this, note that for any $\epsilon > 0$, there exists $T(\epsilon)>0$ such that $\mbox{sup}_{t\ge T(\epsilon)} t\mathbb{P}(d(X_1, \mu) > t) \le \epsilon$ (due to the assumption $n\mathbb{P}(d(X_1, \mu) > n) \underset{n \to \infty}{\to} 0$). Hence, for a fixed $\epsilon > 0$,
\begin{align*}
n^{-1}\int_{0}^{n} t\mathbb{P}(d(X_1, \mu) > t)dt =& \underbrace{n^{-1}\int_{0}^{T(\epsilon)} t\mathbb{P}(d(X_1, \mu) > t)dt}_{\le T(\epsilon)/n} \\ 
&+ \underbrace{n^{-1}\int_{T(\epsilon)}^{n} t\mathbb{P}(d(X_1, \mu) > t)dt}_{\le \epsilon}. 
\end{align*}
Taking $\limsup_{n\to \infty}$ on both sides above, we get 
\[
\limsup_{n\to\infty}n^{-1}\int_{0}^{n} t\mathbb{P}(d(X_1, \mu) > t)dt \le \epsilon.  
\]
Since $\epsilon$ is arbitrary, letting $\epsilon \to 0$ gives $\lim_{n\to\infty}n^{-1}\int_{0}^{n} t\mathbb{P}(d(X_1, \mu) > t)dt = 0$. Therefore, (\ref{eq:int}) goes to zero as $n\to\infty$, as desired.
\end{proof}
 



\section{Technical details for Example 3}

% {\color{red} [I think the computation details (i.e. the rest of example) can be moved to Supp. Material., but should be rewritten.]}
% Let $\mu:=I_{k}$ be a fixed reference point on $\mbox{Sym}^{+}(k)$, where $I_{k}$ is the $k\times k$ identity matrix. Consider a  sequence $(X_i)_{i=1}^{\infty}$ of $\mbox{Sym}^{+}(k)$-valued random variables that are geodesically symmetric about $\mu$ and follow log-normal distributions  with increasing variances. Specifically, for a fixed $\alpha \in (0,1)$, suppose that each independent random variable $X_i$ satisfies $\textsf{vec}\{\mbox{Log}_{\mu}(X_i)\} \sim N_{k(k+1)/2}(\textbf{0}, ~i^{\alpha}I_{k(k+1)/2})$, where $\textsf{vec}(\cdot)$ denotes the vectorization operator \citep[][Section 3.5]{pennec2006riemannian}. 
%
Recall that for each $i\ge 1$, $\textsf{vec}\{\mbox{Log}_{\mu}(X_i)\} \sim N_{k(k+1)/2}(\textbf{0}, ~i^{\alpha}I_{k(k+1)/2})$. 
The Chernoff bound implies that for any $t \in (0, 1/2)$ and for any $x>0$, $\mathbb{P}(\chi^2_{k(k+1)/2} > x) \le e^{-tx}(1-2t)^{-k(k+1)/4}$, where $\chi^2_{k(k+1)/2}$ stands for the chi-squared distribution with degrees of freedom $k(k+1)/2$. In particular, setting $t=1/4$ gives that for any $x > 0$, $\mathbb{P}(\chi^2_{k(k+1)/2} > x) \le 2^{k(k+1)/4}e^{-x/4}$. With this in mind, we get
\begin{align*}
\sum_{i=1}^{n}\mathbb{P}(d(X_i, \mu) > n)&= \sum_{i=1}^{n}\mathbb{P}(\|\mbox{Log}_{\mu}(X_i)\|_{\mu}^2 > n^2) \\
&= \sum_{i=1}^{n}\mathbb{P}(\|\textsf{vec}\{\mbox{Log}_{\mu}(X_i)\}\|_{2}^2 > n^2) \\
&= \sum_{i=1}^{n}\mathbb{P}(\chi^2_{k(k+1)/2}>n^2/i^\alpha) \\
&\le 2^{k(k+1)/4}\sum_{i=1}^{n}\exp(-\{n^2/(4i^\alpha)\}) \\
& < 2^{k(k+1)/4}n\,\exp(-n^{(2-\alpha)}/4) \\
& \underset{n \to \infty}{\to} 0,
\end{align*}
where $\|\cdot \|_{2}$ denotes the $L_2$-norm in $\mathbb{R}^{k(k+1)/2}$. In addition, 
\begin{align*}
n^{-2}\sum_{i=1}^{n}\mathbb{E}[d^2(\overline{X_{n,i}},\mu)] 
&\le n^{-2}\sum_{i=1}^{n}\mathbb{E}[d^2(X_i,\mu)] = n^{-2}\sum_{i=1}^{n}\mathbb{E}[\|\mbox{Log}_{\mu}(X_i)\|_{\mu}^2] \\
&=n^{-2}\sum_{i=1}^{n}\textsf{tr}[\textsf{Cov}\{\textsf{vec}(\mbox{Log}_{\mu}(X_i))\}] = \frac{k(k+1)}{2}n^{-2}\sum_{i=1}^{n}i^\alpha \\
&\le \frac{k(k+1)}{2} n^{-2}\int_{1}^{n+1}x^\alpha dx \\
&= \frac{k(k+1)}{2}n^{-2}\{(n+1)^{\alpha+1} -1\}/(\alpha+1) \\
&< \frac{k(k+1)}{2}n^{\alpha-1}(1+ 1/n)^{\alpha + 1}/(\alpha+1) \\
&\underset{n \to \infty}{\to} 0.
\end{align*}
Since the hypotheses of Theorem 1(i) hold, $\mu_n(X_1, X_2, \ldots , X_n) \overset{\mathbb{P}}{\to} \mu$ as $n\to \infty$ by Theorem 1(i). 






\section{A partial converse for Theorem 1(ii) and its proof.}
 

An interesting question is whether the converse of Theorem 1 for case (ii) holds. This concerns the extent to which the Kolmogorov--Feller tail condition is tight for the weak law. In the Euclidean space the following result is classical; it follows from
an application of Corollary 3.3 in \cite{marcus1979stable}. 
For completeness, we include a proof in this section. 

\begin{theorem}\label{thm:main2}
 Suppose that $M=\mathbb{R}^k$ is endowed with an arbitrary inner product, and that Condition (C) holds and the $X_i$ are identically distributed. If $\mu_n \overset{\mathbb{P}}{\to} \mu$ as $n \to \infty$, then $n\mathbb{P}(\|X_1-\mu\| > n) \underset{n \to \infty}{\to} 0$.
\end{theorem}
Since a noncompact symmetric space with identically zero sectional curvature is isometric to a Euclidean space, the result established in the Euclidean setting applies verbatim to this flat symmetric-space case. 
We conjecture that the result of Theorem \ref{thm:main2} does not hold for general non-compact symmetric spaces.

To verify Theorem~\ref{thm:main2}, we will use the following lemma regarding a symmetric distribution on the real line. 
\begin{lemma}\label{lem:same_dist}
Given a real-valued random variable $Y$, suppose that $Y$ and $-Y$ have the same distribution; that is, $Y\overset{d}{=}-Y$. Then, $\mathbb{P}(Y\ge 0) \ge 1/2$. 
\end{lemma}

\begin{proof}[Proof of Lemma~\ref{lem:same_dist}]
Since $Y$ and $-Y$ have the same distribution, an application of countable sub-additivity of probability measures yields that 
\begin{align*}
1=\mathbb{P}(Y\in \mathbb{R}) &\le \mathbb{P}(Y \ge 0) + \mathbb{P}(Y\le 0) \\
&= \mathbb{P}(Y \ge 0) + \mathbb{P}(-Y\le 0) \\
&\le 2\mathbb{P}(Y\ge 0).
\end{align*}
In turn, this leads to the desired result.
\end{proof} 
 


\begin{proof}[Proof of Theorem~\ref{thm:main2}]
When $k=1$, it is well-known that the statement of Theorem~\ref{thm:main2} holds; see, for example, \cite{stoica2010kolmogorov, yuan2015conditional, naderi2019version}. For this reason, suppose that $M$ is equal to $\mathbb{R}^k$, $k\ge 2$, endowed with an arbitrary inner product $\langle~, ~\rangle$ and the induced norm $\|\mathbf{x}\|:=\langle \mathbf{x}, \mathbf{x}\rangle^{1/2}$ for $\mathbf{x}\in \mathbb{R}^k$.

Without loss of generality, we may assume that $\mu=\textbf{0}\in \mathbb{R}^k$. (If not, for any $1\le i\le n$, replace $X_i$ with $X_i - \mu$.) Suppose, for contradiction, that there exists a $\delta > 0$ and an increasing subsequence of natural numbers $(n_{j})_{j\ge 1}$ such that for any $j\ge 1$, $n_j\mathbb{P}(\|X_1\| > n_j) \ge \delta$. Notably, $\mathbb{P}(\|X_1\|>n)$ is close to zero for large enough $n$ and an inequality $(1-p)^n \le e^{-np}$ holds for small $p\ge0$. With these in mind, for large enough $j$
\begin{align}\label{ineq:strict}
\mathbb{P}(\exists 1 \le i\le n_{j}: \|X_i\|> n_{j}) &= 1- \mathbb{P}(\max_{1\le i\le n_{j}} \|X_i\| \le n_{j}) \nonumber \\
&= 1-\{\mathbb{P}(\|X_1\| \le n_{j})\}^{n_{j}} \nonumber \\
&= [1-\{1-\mathbb{P}(\| X_1\|>n_{j})\}]^{n_{j}} \nonumber \\
&\ge 1-e^{-n_{j}\mathbb{P}(\|X_1 \|>n_{j})} \nonumber \\
&\ge 1-e^{-\delta} ~(>0).
\end{align}
Let $j$ be a fixed large natural number and denote $S_{n_{j}}=\sum_{\ell=1}^{n_{j}}X_\ell$ and $S_{-i} = S_{n_{j}} - X_i ~ (1\le i\le n_j)$. An application of (\ref{ineq:strict}) yields that
\begin{align}\label{ineq:strict2}
\mathbb{P}(\|S_{n_j}/n_{j}\|>1) \ge& \mathbb{P}(\exists1\le i\le n_j: \|X_i\|>n_j ~ \mbox{and} ~ \langle X_i, S_{-i} \rangle \ge 0) \nonumber \\
=& \mathbb{P}(\langle X_i, S_{-i}\rangle \ge 0  \mid \exists1\le i\le n_j: \|X_i\|>n_j) \nonumber \\
& \times \mathbb{P}(\exists1\le i\le n_j: \|X_i\|>n_j) \nonumber \\
\ge& 1/2 \times \mathbb{P}(\exists 1\le i\le n_j: \|X_i\|>n_j) \nonumber \\
\ge& (1-e^{-\delta})/2 ~( > 0),
\end{align}
where the first inequality holds due to the inequality below:
$$
\|S_{n_{j}}\|^2=\langle X_i + S_{-i}, X_i+S_{-i} \rangle =\|X_i\|^2 + 2\langle X_i, S_{-i} \rangle +\|S_{-i} \|^2 > n_j^2.
$$
Moreover, the second inequality of (\ref{ineq:strict2}) is satisfied due to an application of Lemma~\ref{lem:same_dist} to the fact that $\langle X_i, S_{-i}\rangle \overset{d}{=} \langle X_i, -S_{-i} \rangle$ (recall that $\mathbb{P}_{X}$ is geodesically symmetric about $\mathbf{0}$, thereby implying $S_{-i}\overset{d}{=}-S_{-i}$). Since $\mu_n=S_n/n$ and (\ref{ineq:strict2}) holds for any sufficiently large $j$, $\mu_n$ does not converge to $\mu ~(=\mathbf{0})$ in probability as $n\to \infty$. This contradicts the assumption.
\end{proof}


\section{Sharpness of Theorem 1(ii)} %Kolmogorov--Feller tail condition}

The Kolmogorov--Feller tail condition, namely $n\mathbb{P}(d(X_1, \mu_0) > n) \to 0$, implies that $\mathbb{E}[d^{\alpha}(X, \mu)] <\infty$ for all $\alpha \in (0, 1)$. Therefore, under the same conditions as in Theorem 1(ii), the $\alpha$-\Frechet\ mean 
$$\mu_n^{\alpha}:=\arg\min_{m \in M} n^{-1}\sum_{i=1}^{n}d^\alpha(X_i, m)$$ 
for any $\alpha \in (0,1)$, converges to $\mu$ almost surely as $n$ increases \cite[][Corollary 5.2]{schotz2022strong}. Nevertheless, Theorem 1(ii) is sharp in the sense that, under the same Kolmogorov--Feller tail condition, the convergence in probability established in Theorem 1 is not, in general, strengthened to almost sure convergence. The following example demonstrates this.

\begin{example}[Failure of almost sure convergence]
Let $(M, d)$ be a non-compact symmetric space, fix $\mu \in M$, and let $\gamma:\mathbb{R}\to M$ be a unit-speed geodesic with $\gamma(0)=\mu$. Let $Z$ be a symmetric real-valued random variable satisfying
$$
\mathbb{P}(Z = \pm 2^i) = \frac{1}{(2i)2^i} \quad (i \ge 1), \quad \mbox{and} \quad \mathbb{P}(Z=0) = 1 - \sum_{i=1}^\infty \frac{1}{i(2^i)}.
$$
Define the $M$-valued random variable $X= \gamma(Z)$, and let $X, X_1, X_2, \ldots$ be independent and identically distributed random variables. Then, $\mathbb{P}_X$ is geodesically symmetric about $\mu$. Since $\gamma$ is unit-speed, $d(X,\mu) = |Z|$. The Kolmogorov--Feller condition is satisfied:
$
n\mathbb{P}(d(X,\mu)>n) = n\mathbb{P}(|Z|>n) \to 0
$
as $n\to \infty$, which in turn implies that $\mu_n \to \mu$ in probability. On the other hand, since all observations lie on the geodesic $\gamma(\mathbb{R})$, the sample Fréchet mean $\mu_n=\mu_n(X_1, X_2, \ldots, X_n)$ coincides with the Euclidean sample mean along that geodesic, that is,
$$
\mu_n = \gamma(\bar Z_n),
\qquad
\bar Z_n := \frac{1}{n}\sum_{i=1}^n Z_i,
$$
where $Z, Z_1, Z_2, \ldots$ are independent and identically distributed random variables. However, $\bar Z_n \not\to 0$ almost surely. Thus,
$$
\mu_n \not\to \mu \qquad \text{almost surely}.
$$
Therefore, the conditions in Theorem 1(ii) do not, in general, guarantee almost sure convergence.
\end{example}
%Recall that $Z$ is a symmetric real-valued random variable satisfying
%$$
%\mathbb{P}(Z = \pm 2^i) = \frac{1}{(2i)2^i} \quad (i \ge 1), \quad \mbox{and} \quad \mathbb{P}(Z=0) = 1 - \sum_{i=1}^\infty \frac{1}{i(2^i)}.
%$$
%Recall also that $X= \gamma(Z)$ and $d(X,\mu) = |Z|$. 
Technical details for Example E.1 are described below. For any integer $i\ge 1$,
\begin{align*}
2^i \mathbb{P}(|Z| > 2^i) &= \sum_{m=1}^{\infty} \frac{1}{i+m}\cdot\frac{1}{2^m} \\
&< \frac{1}{i+1} (\sum_{m=1}^{\infty} \frac{1}{2^m}) \\
&= \frac{1}{i+1} \\
&\underset{i \to \infty}{\to} 0.
\end{align*}
For each $n\ge 2$, there exists a unique integer $i\ge 1$ such that $2^i \le n <2^{i+1}$. Then,
$$
n\mathbb{P}(|Z| > n) <2^{i+1}\mathbb{P}(|Z|>2^i) = 2 \cdot 2^{i}\mathbb{P}(|Z|>2^i). 
$$
As $n$ increases, so does $i$. Combining these facts, we have
$$
n\mathbb{P}(d(X,\mu)>n) = n\mathbb{P}(|Z|>n) \underset{n\to \infty}{\to} 0.
$$
So the Kolmogorov--Feller tail condition is fulfilled. %Since all observations lie on the geodesic $\gamma(\mathbb{R})$, the sample Fréchet mean $\mu_n$ coincides with the Euclidean sample mean along that geodesic, that is,
Also, recall that 
$$
\mu_n = \gamma(\bar Z_n),
\qquad
\bar Z_n := \frac{1}{n}\sum_{i=1}^n Z_i,
$$
where $Z, Z_1, Z_2, \ldots$ are independent and identically distributed random variables. 

Now, we wish to verify that $\sum_{n=1}^{\infty}\mathbb{P}(|Z_n| > n)=\infty$. By Cauchy's condensation test, 
$$
\sum_{n=1}^{\infty}\mathbb{P}(|Z_n| > n)= \sum_{n=1}^{\infty}\mathbb{P}(|Z| > n)=\infty \quad \Leftrightarrow \quad \sum_{n=1}^{\infty} 2^n \mathbb{P}(|Z|>2^n) = \infty.
$$
Since $2^n\mathbb{P}(|Z|>2^n)= \sum_{m=1}^{\infty} \frac{1}{2^m(n+m)} > \frac{1}{2(n+1)}$ and $\sum_{n=1}^{\infty}\frac{1}{2(n+1)}=\infty$, $$\sum_{n=1}^{\infty}\mathbb{P}(|Z_n| > n) =\infty.$$ 
By independence on $(Z_n)_{n\ge 1}$ and the second Borel--Cantelli lemma,
$$
\mathbb{P}(\underbrace{|Z_n|>n \quad \text{infinitely many} ~ n}_{=:E}) = 1,
$$ 
On the event $E$, there exists a subsequence $(n_k)_{k\ge 1}$ such that $\frac{|Z_{n_k}|}{n_k} > 1$, and hence $Z_n/n \not\to 0$. In particular, $\bar Z_n \not\to 0$ whenever $E$ occurs. Consequently,
$$
\mu_n \not\to \mu,
$$
almost surely.



\end{appendix}
%Use the standard \LaTeX\ commands for headings in \verb|{appendix}|.
%Headings and other objects will be numbered %automatically.
%\begin{equation}
%\mathcal{P}=(j_{k,1},j_{k,2},\dots,j_{k,m(k)}). \label{path}
%\end{equation}

%Sample of cross-reference to the formula (\ref{path}) in Appendix \ref{appB}.



% \begin{supplement}
% \stitle{Supplement to ``Huber means on Riemannian manifolds"}
% \sdescription{This supplementary material provides the auxiliary statements for additional theoretical properties of the Huber mean, technical details, additional lemmas and propositions, and all proofs of findings.}
% \end{supplement}

%%%%%%%%%%%%%%%%%%%%%%%%%%%%%%%%%%%%%%%%%%%%%%%%%%%%%%%%%%%%%
%%                  The Bibliography                       %%
%%                                                         %%
%%  imsart-???.bst  will be used to                        %%
%%  create a .BBL file for submission.                     %%
%%                                                         %%
%%  Note that the displayed Bibliography will not          %%
%%  necessarily be rendered by Latex exactly as specified  %%
%%  in the online Instructions for Authors.                %%
%%                                                         %%
%%  MR numbers will be added by VTeX.                      %%
%%                                                         %%
%%  Use \cite{...} to cite references in text.             %%
%%                                                         %%
%%%%%%%%%%%%%%%%%%%%%%%%%%%%%%%%%%%%%%%%%%%%%%%%%%%%%%%%%%%%%

%% if your bibliography is in bibtex format, uncomment commands:
%\bibliographystyle{unsrt}  
%\bibliography{hmr}  

\bibliographystyle{imsart-nameyear}
\bibliography{KFLLN}     

\end{document}
